% from ega4-2-fr.tex
\oldpage[IV-2]{164}

\item[b)] Pour tout sous-préschéma fermé intègre $Y$ de $X$,
$\operatorname{Reg}(Y)$ contient une partie ouverte non vide de $Y$.
\end{enumerate}
\end{corollary}

\begin{proof}
Il est clair que a) entraîne b), car si $Y$ est intègre et $y$ est son point
générique, $\mathcal O_{Y,y}$ est un corps, donc $y\in\operatorname{Reg}(Y)$.
Inversement, pour voir que b) entraîne a), considérons un sous-préschéma
fermé intègre $Z$ de $Y$ de point générique $z$; si $Y'$ est le
sous-préschéma intègre de $X$ ayant pour espace sous-jacent l'adhérence
$\overline{\{z\}}$ de $z$ dans $X$, $Z$ est ouvert dans $Y'$ et le
sous-préschéma $Z$ de $Y'$, étant réduit, est induit par $Y'$ sur l'ouvert
$Z$ de l'espace sous-jacent à $Y'$. Or l'hypothèse entraîne que
$\operatorname{Reg}(Y')$ est un voisinage de $z$ dans $Y'$, donc
$\operatorname{Reg}(Z)$ est un voisinage de $z$ dans $Z$; il suffit alors
d'appliquer (6.12.2) en y remplaçant $X$ par $Y$ et $Y$ par $Z$.
\end{proof}

\begin{theorem}[6.12.4]
\label{IV.2.s6.12.4-fr}
(Nagata). Soient $A$ un anneau noethérien,
$X=\operatorname{Spec}(A)$. Les conditions suivantes sont équivalentes:
\begin{enumerate}
\item[a)] Pour tout préschéma $X'$ localement de type fini sur $X$,
$\operatorname{Reg}(X')$ est ouvert dans $X'$.
\item[b)] Pour toute $A$-algèbre $A'$ finie et intègre, il existe une
partie ouverte non vide de $X'=\operatorname{Spec}(A')$ contenue dans
$\operatorname{Reg}(X')$.
\item[c)] Pour tout idéal premier $\mathfrak p$ de $A$ et toute extension
finie radicielle $K'$ du corps des fractions $K$ de $A/\mathfrak p$, il
existe une sous-$A$-algèbre $A'$ de $K'$, finie sur $A$, ayant $K'$ pour
corps des fractions, et telle qu'il existe dans
$X'=\operatorname{Spec}(A')$ un ensemble ouvert non vide contenu dans
$\operatorname{Reg}(X')$.
\end{enumerate}
\end{theorem}

\begin{proof}
Il est clair que a) implique b). Pour voir que b) entraîne c), il suffit de
remarquer qu'on peut prendre pour générateurs de l'extension $K'$ de $K$
des éléments entiers sur $A/\mathfrak p$ (et a fortiori sur $A$), et comme
ces éléments sont en nombre fini, ils engendrent une $A$-algèbre finie
$A'$ dont $K'$ est le corps des fractions; on peut alors appliquer b) à
$A'$. Reste donc à prouver que c) entraîne a). La question étant locale sur
$X'$, on peut d'abord supposer que $X'=\operatorname{Spec}(A')$, où $A'$
est une $A$-algèbre de type fini; compte tenu de (6.12.2), il suffit de
prouver que pour tout sous-préschéma fermé intègre $Y'$ de $X'$,
$\operatorname{Reg}(Y')$ contient une partie ouverte non vide de $Y'$.
Autrement dit, on peut se borner à prouver que si $A'$ est une $A$-algèbre
intègre de type fini et $X'=\operatorname{Spec}(A')$,
$\operatorname{Reg}(X')$ contient une partie ouverte non vide de $X'$.
Soit $K'$ le corps des fractions de $A'$; si $\mathfrak p$ est l'image
canonique dans $X$ du point générique de $X'$, $K'$ est une extension du
corps des fractions $K$ de $A/\mathfrak p$, et $A/\mathfrak p$
s'identifie à un sous-anneau de $A'$, $A'$ étant une
$(A/\mathfrak p)$-algèbre de type fini. On peut évidemment se borner dans
la suite de la démonstration au cas $\mathfrak p=0$. Distinguons maintenant
deux cas:

I) $K'$ est une extension séparable de $K$. --- On est alors ramené à
prouver le

\begin{lemma}[6.12.4.1]
\label{IV.2.s6.12.4.1-fr}
Soient $A$ un anneau intègre noethérien, $A'$ une $A$-algèbre intègre de
type fini, contenant $A$ et telle que le corps des fractions $K'$ de $A'$
soit une extension séparable du corps des fractions $K$ de $A$. Si
$\operatorname{Spec}(A)$ contient un ouvert non vide formé de points
réguliers, il en est de même de $\operatorname{Spec}(A')$.
\end{lemma}

\begin{proof}
Remplaçant $A$ par un anneau de fractions $A_f$ de sorte que
$D(f)\subset\operatorname{Reg}(\operatorname{Spec}(A))$, on peut déjà
supposer l'anneau $A$ régulier. Par hypothèse il existe dans $K'$ un système
$(t_i)_{1\leq i\leq n}$ d'éléments algébriquement indépendants sur $K$ et
tels que $K'$ soit une extension algébrique séparable finie de
$K(t_1,\ldots,t_n)$; en considérant un système fini
